% ======================================================================
%  Section 4.3.2 --- Robust normalization of the residuals.  Corrected.
%
%  VERIFIED AGAINST THE IMPLEMENTATION. ViSP computes
%      m_mad = 1.4826 * normmedian;
%  and enforces m_mad_min = 0.0017. Both numbers in the text are right.
%
%  ONE BROKEN PROMISE, ONE CLAIM NOW MEASURABLE, THREE SMALL THINGS.
%
%  A. "The admissible occlusion fraction ... is characterized empirically in
%     Section 4.7.2." It is not. Section 4.7.2 tests a single occluder at a
%     single size. There is no sweep, and the draft of 4.7.2 explicitly
%     lists an occluded fraction approaching the breakdown point as future
%     work. The promise has to go or be redirected.
%
%  B. The dilation argument, by contrast, CAN now be supported, because it
%     is visible in the iteration-600 weight map. Measured on that map:
%
%       occluder, the connected region where G vanishes   3829   5.80%
%       largest connected rejected region (w < 0.25)      5886   8.92%
%       ratio                                             1.54
%       equivalent radii                          34.9 px -> 43.3 px
%       i.e. a rejected ring of about 8 px around the disc
%       against a nominal filter half-width of 4 px (9x9 kernels)
%
%     So the prediction holds and is quantified. The section should predict
%     and point forward; 4.7.3 should report the number.
%
%  C. NOTE FOR 4.7.2 AND 4.7.3. The occluder covers 7.4% of the
%     measurements at ITERATION 1 and 5.80% at ITERATION 600, because the
%     camera starts at Z = 1.10 m and finishes at Z = 1.30 m, so the disc
%     subtends a smaller angle at the end. The ratio of areas predicted by
%     the depths alone is (1.30/1.10)^2 = 1.40 against 7.4/5.8 = 1.28
%     measured, the difference being the threshold on G. Both figures are
%     correct at their own iteration, but 4.7.3 currently compares the 9.5%
%     rejected at iteration 600 against the 7.4% measured at iteration 1.
%     That comparison must use 5.80%, which makes the point more strongly,
%     not less: 9.5% against 5.80% is a clear excess, where 9.5% against
%     7.4% looked like the occluder plus noise.
% ======================================================================

\subsection{Robust normalisation of the residuals}
\label{sec:mad}
%%% Sections 4.3.1, 4.7.2 and 4.7.3 all refer here; the label is needed.

The normalisation $u_{j}=e_{j}/\hat{s}$ requires a scale estimate $\hat{s}$
that is itself insensitive to the outliers it is meant to detect; the sample
standard deviation is unsuitable, since the occluded measurements inflate it
and thereby mask their own detection. The median absolute deviation is used
instead,
\begin{equation}
    \mathrm{MAD}=\operatorname{med}_{j}\bigl\lvert e_{j}-\operatorname{med}_{k}(e_{k})\bigr\rvert,
    \qquad
    \hat{s}=1.4826\,\mathrm{MAD},
    \label{eq:mad}
\end{equation}
the constant $1.4826$ making $\hat{s}$ a consistent estimator of the standard
deviation for normally distributed residuals, with the same proviso as
attaches to the tuning constant of Section~\ref{sec:tukey-biweight}: the
normalisation is adopted as the conventional one, not because the residuals
of a direct scheme satisfy the model from which it is derived.

\paragraph{Breakdown point and occluded area.}
The MAD has a breakdown point of $50\%$, so the scale estimate remains
meaningful provided that fewer than half of the residual entries are
contaminated. This does not translate into an equal bound on the occluded
image area. Because each Hermite response draws on a neighbourhood determined
by the effective filter support, and because the four scales are stacked, an
occluded region can affect responses extending beyond the directly occluded
pixels: the set of \emph{potentially affected} responses is approximately the
occluded region dilated by that support. The ratio of affected responses to
occluded pixels therefore grows with the perimeter-to-area ratio of the
occlusion, and is larger for a small or fragmented mask than for a single
compact one. The admissible occlusion fraction is consequently smaller than
the nominal breakdown point.

%%% (B) The prediction is now testable and tested. Note the careful
%%% distinction: dilation is a statement about which responses are AFFECTED,
%%% and only at convergence does the set of REJECTED responses come to
%%% resemble it. At iteration 1 the rejected set is smaller than the
%%% occluder, as 4.7.3 reports, because the robust scale is still large.
This dilation is directly observable in the weight maps of
Section~\ref{sec:weight-maps}. At convergence the rejected region forms a
ring around the occluding disc some two filter half-widths wide, and exceeds
the occluded region itself by a factor of about $1.5$ --- the measurement is
reported there. It should be noted that the two sets coincide only at
convergence: while the pose error is still large the robust scale is large in
proportion and the rejected set is \emph{smaller} than the occluded one, so
the dilation bounds what the occlusion can affect rather than what the
estimator rejects at any given instant.

%%% (A) The promise of an empirical characterisation is replaced by an
%%% honest statement of what was and was not done. A single occluded
%%% fraction was tested; the bound itself was not located.
The experiments of Section~\ref{sec:ablation-results} use a single compact
occluder covering a fraction well below this bound, and therefore do not
locate it. Determining the admissible fraction, and in particular the
dependence on mask fragmentation that the argument above predicts, would
require a sweep over occluder size and shape, which is not attempted here.

\paragraph{Implementation.}
The median absolute deviation in~\eqref{eq:mad} is lower-bounded by the
minimum value $0.0017$ enforced by the robust estimator, which prevents
instability when the residual dispersion approaches zero. In the runs reported
in Section~\ref{sec:results} the dispersion remains several orders of
magnitude above this bound throughout, so the floor is never active; it
guards against a degenerate case rather than shaping the behaviour observed.

The resulting weights are
\begin{equation}
    w^{\mathrm{T}}_{j}=w^{\mathrm{T}}\!\left(\frac{e_{j}}{\hat{s}}\right),
    \qquad j=1,\dots,P,
    \label{eq:tukey-weights}
\end{equation}
recomputed at every servoing iteration from the current residuals, and
assembled into
$\mathbf{D}^{\mathrm{T}}=\mathrm{diag}(w^{\mathrm{T}}_{1},\dots,w^{\mathrm{T}}_{P})$.

% ======================================================================
%  D. A NOTATION CLASH WORTH FIXING BEFORE IT SPREADS
%
%  D^T here means "the D of the Tukey variant", while \top is transpose
%  throughout 4.2.2 -- and (4.4) contains L_R^\top D^2 L_R, so the reader
%  meets both within a few lines. Section 4.4.3 then adds D^H. Consider
%
%      D_T and D_H   (subscripts, no collision with transpose)
%  or  D^{(2)} and D^{(3)}   (by variant number, which also ties the
%                             notation to the ablation of 4.6.1)
%
%  Mathematically harmless, since D is symmetric, but the reader has to
%  decide each time which superscript is which.
%
%  E. SMALLER
%  - "normalization" / "characterized" are -ize; the chapter is -ise
%    elsewhere. The heading is corrected above; check the rest.
%  - Section 4.7.2 and Table 4.2's caption both state the occluder as 7.4%
%    of the measurements without saying at which pose. See note C above.
% ======================================================================
